% Outline for the land-model section of the VFS paper.

\section{Land surface model}\label{sec:land}

\subsection{Design goals and state ownership}
\begin{itemize}
  \item Motivate a tensor-native, dense land model suitable for shared-cell coupling and accelerator execution.
  \item Define the land state: layered liquid soil water, layered soil temperature, optional soil ice, canopy water, snow-water equivalent, skin temperature, and static land-use type.
  \item State units: water stores are kg m$^{-2}$, temperatures are K, and soil fields use shape $(n_{\mathrm{cell}}, n_{\mathrm{soil}})$.
  \item Explain that the surface system blends land boundary properties and routes transfers, while the land component owns its prognostic stores.
\end{itemize}
\subsection{Hydrology and vegetation controls}
\begin{itemize}
  \item Describe layered bucket hydrology: precipitation and snowmelt enter the upper layer, saturation excess cascades downward, and gravity drainage removes lower-boundary water.
  \item Describe root-weighted evaporation extraction above the wilting point and moisture availability between wilting and field capacity.
  \item Describe static land-use lookup fields: albedo, roughness, evaporation factor, root depth, emissions factor, and canopy capacity factor.
  \item Describe canopy interception, throughfall, canopy evaporation, runoff, and drainage in terms of explicit water transfers.
\end{itemize}

\subsection{Snow and energy balance}
\begin{itemize}
  \item Describe snow-water storage, fractional snow cover, snow-modified albedo and roughness, and snow insulation of the soil upper boundary.
  \item Give the local skin-energy balance with net radiation, turbulent fluxes, conductive soil flux, and energy-limited snowmelt.
  \item Describe explicit soil thermal diffusion and its stability substepping.
  \item Define sign conventions for downward radiative fluxes and fluxes leaving the skin reservoir.
\end{itemize}

\subsection{Coupling and conservation accounting}
\begin{itemize}
  \item Explain the \texttt{ConservativeTransfer} ledger: atmosphere-to-land rain/snow, land-to-atmosphere evaporation, and land-to-runoff export.
  \item Explain fractional land masks and the requirement that inactive cells retain their land state.
  \item Define total stored land water as soil liquid plus soil ice plus snow plus canopy water, area weighted over land fraction.
\end{itemize}

\subsection{Verification and results to report}
\begin{table}[t]
\centering
\caption{Planned land-model verification cases.}
\label{tab:land-verification}
\begin{tabular}{p{0.28\linewidth}p{0.35\linewidth}p{0.27\linewidth}}
\hline
Case & Required diagnostic & Planned paper result \\
\hline
Synthetic water budget & Storage change versus transfer ledger; runoff and drainage closure & Residual table at float64 and float32 \\
Diurnal rain/snow forcing & Soil water, canopy water, snow, runoff, skin temperature & Multi-day time series and budget residual \\
Snowmelt energy limit & Potential versus energy-limited melt; skin energy residual & Process-level figure \\
Land-use contrast & Albedo, roughness, root profile, evaporation, and runoff across classes & Comparison table or panel plot \\
Restart/mask test & Restart equality and inactive-cell invariance & Error norm summary \\
\hline
\end{tabular}
\end{table}

\subsection{Known scope limits}
\begin{itemize}
  \item Describe the implementation as a layered, conservation-aware prototype rather than a production terrestrial biosphere model.
  \item State absent processes: lateral groundwater and river routing, dynamic vegetation, carbon and nutrient cycles, detailed soil phase change, and calibrated global land datasets.
  \item State that long coupled restart and climate validation remain required before making climate-performance claims.
\end{itemize}
